\documentclass[11pt,a4paper,openany]{ctexbook}

\usepackage[margin=2.2cm]{geometry}
\usepackage{enumitem}
\usepackage{booktabs}
\usepackage{xcolor}
\usepackage{hyperref}
\usepackage{tcolorbox}
\usepackage{listings}
\usepackage{fancyhdr}
\usepackage{titlesec}

\definecolor{codebg}{RGB}{248,248,248}
\definecolor{notebg}{RGB}{240,248,255}
\definecolor{warnbg}{RGB}{255,248,240}

\lstset{
  basicstyle=\ttfamily\small,
  backgroundcolor=\color{codebg},
  frame=single,
  breaklines=true,
  columns=fullflexible,
  keepspaces=true,
  showstringspaces=false,
  extendedchars=false
}

\lstdefinelanguage{yaml}{}

\newtcolorbox{note}[1][]{
  colback=notebg, colframe=blue!40!black, title=提示, #1
}
\newtcolorbox{warn}[1][]{
  colback=warnbg, colframe=orange!70!black, title=注意, #1
}

\pagestyle{fancy}
\fancyhf{}
\fancyhead[LE,RO]{\thepage}
\fancyhead[LO]{\nouppercase{\rightmark}}
\fancyhead[RE]{\nouppercase{\leftmark}}
\renewcommand{\headrulewidth}{0.4pt}

\title{Elasticsearch 8.19.23 安装与排错记录\\\large Ubuntu / deb 包部署}
\author{}
\date{}

\begin{document}
\frontmatter
\maketitle

\begin{note}
本文记录在 Ubuntu 上通过 \texttt{.deb} 包安装 Elasticsearch 8.19.23 的完整过程，以及安装后遇到的安全认证失败、集群 RED、磁盘水位线、数据目录迁移等问题的排查与解决方法。不含数据迁移内容。
\end{note}

\tableofcontents
\mainmatter

\chapter{安装前检查}

安装前先确认磁盘空间与目录规划，可以避免安装后因磁盘水位线导致的集群 RED、安全索引初始化失败等连锁问题。本次排错中，磁盘可用空间仅 5.5\% 是多个故障的根因之一。

\section{检查可用磁盘空间}

查看根分区及计划用作数据目录的分区：

\begin{lstlisting}[language=bash]
df -h /
df -h /data
\end{lstlisting}

Elasticsearch 默认磁盘水位线规则：
\begin{itemize}[leftmargin=1.4em]
  \item 低水位线 85\%：不再向该节点分配新分片；
  \item 高水位线 90\%：尝试将该节点上的分片迁走；
  \item 洪水水位线 95\%：每个索引只保留一个分片。
\end{itemize}

\begin{warn}
数据目录所在分区可用空间应保持在 80\% 以下。若根分区紧张，应在安装前规划独立的 \texttt{path.data} 分区，而不是安装后再迁移。
\end{warn}

\section{规划数据目录}

默认数据目录为 \texttt{/var/lib/elasticsearch}，位于根分区。建议在安装前确定数据目录位置：

\begin{lstlisting}
/var/lib/elasticsearch      # 默认，位于根分区
/data/elasticsearch         # 独立分区示例
\end{lstlisting}

若使用独立分区，提前创建目录并设置权限：

\begin{lstlisting}[language=bash]
sudo mkdir -p /data/elasticsearch
sudo chown elasticsearch:elasticsearch /data/elasticsearch
sudo chmod 750 /data/elasticsearch
\end{lstlisting}

\begin{note}
\texttt{elasticsearch} 用户在安装包解压时才会创建。若目录在安装前创建，可先占位，安装后再修正属主。
\end{note}

\section{检查内存与文件描述符}

\begin{lstlisting}[language=bash]
free -h
ulimit -n
\end{lstlisting}

Elasticsearch 默认堆内存为物理内存的一半，且要求 \texttt{vm.max\_map\_count} 不低于 262144。deb 包通常会写入相关系统配置，安装后可复查：

\begin{lstlisting}[language=bash]
sysctl vm.max_map_count
\end{lstlisting}

\chapter{安装与启动}

\section{安装方式}

使用 \texttt{dpkg} 安装官方 \texttt{.deb} 包：

\begin{lstlisting}[language=bash]
dpkg -i /root/Packages/elasticsearch-8.19.23-amd64.deb
\end{lstlisting}

安装过程中会自动创建 \texttt{elasticsearch} 用户和组，并完成安全自动配置（Security autoconfiguration）。终端会输出：

\begin{itemize}[leftmargin=1.4em]
  \item \texttt{elastic} 内置超级用户的初始密码；
  \item TLS 已为 transport 和 HTTP 层启用并配置；
  \item Kibana / 节点 enrollment token 的生成方式。
\end{itemize}

\begin{warn}
首次安装输出的 \texttt{elastic} 初始密码务必立即保存。该密码只在此刻显示一次。
\end{warn}

\section{启动服务}

\begin{lstlisting}[language=bash]
sudo systemctl daemon-reload
sudo systemctl enable elasticsearch.service
sudo systemctl start elasticsearch.service
\end{lstlisting}

\section{证书与配置文件位置}

deb 包安装方式下，HTTP 层 CA 证书位于：

\begin{lstlisting}
/etc/elasticsearch/certs/http_ca.crt
\end{lstlisting}

配置文件位于 \texttt{/etc/elasticsearch/elasticsearch.yml}。

\chapter{基础配置}

\section{单节点模式}

在单机部署场景中，必须显式设置单节点发现模式，并\textbf{移除}多节点首次选举配置。

\begin{lstlisting}[language=yaml]
discovery.type: single-node
\end{lstlisting}

\begin{warn}
\texttt{discovery.type: single-node} 与 \texttt{cluster.initial\_master\_nodes} 互斥。后者必须删除或注释掉，否则节点会以 \texttt{multi-node} 模式启动，导致集群形成过程异常。
\end{warn}

启动日志中可确认实际生效的发现模式：

\begin{lstlisting}
using discovery type [single-node] and seed hosts providers [settings]
\end{lstlisting}

若日志显示 \texttt{multi-node}，说明单节点配置未生效。

\section{数据目录}

默认数据目录为 \texttt{/var/lib/elasticsearch}。可通过 \texttt{path.data} 指向其他分区：

\begin{lstlisting}[language=yaml]
path.data: /data/elasticsearch
\end{lstlisting}

修改后需确保目录属主与权限正确：

\begin{lstlisting}[language=bash]
sudo mkdir -p /data/elasticsearch
sudo chown elasticsearch:elasticsearch /data/elasticsearch
sudo chmod 750 /data/elasticsearch
\end{lstlisting}

\chapter{认证与密码操作}

\section{验证认证}

使用 CA 证书和 \texttt{elastic} 用户访问：

\begin{lstlisting}[language=bash]
curl --cacert /etc/elasticsearch/certs/http_ca.crt \
  -u elastic:'<密码>' https://localhost:9200
\end{lstlisting}

返回 JSON 中 \texttt{version.number} 为 \texttt{8.19.23} 即表示认证成功。

\section{重置密码}

官方工具：

\begin{lstlisting}[language=bash]
sudo /usr/share/elasticsearch/bin/elasticsearch-reset-password -u elastic -i
\end{lstlisting}

\begin{warn}
该工具需要集群健康状态不是 RED 才能正常工作。若集群为 RED，会报 \texttt{exit code 69}。
\end{warn}

\section{file realm 应急用户}

当 \texttt{elastic} 用户无法认证、且重置工具不可用时，可通过 file realm 创建临时超级用户绕过：

\begin{lstlisting}[language=bash]
# 停止服务后创建
sudo systemctl stop elasticsearch
sudo /usr/share/elasticsearch/bin/elasticsearch-users \
  useradd recovery_admin -p 'TempRecovery123!' -r superuser
sudo systemctl start elasticsearch
\end{lstlisting}

file realm 用户信息存储在配置目录下的 \texttt{users} 与 \texttt{users\_roles} 文件中，不依赖 \texttt{.security} 索引，因此可在安全索引损坏时提供登录入口。

用临时用户重置 \texttt{elastic} 密码：

\begin{lstlisting}[language=bash]
curl --cacert /etc/elasticsearch/certs/http_ca.crt \
  -u recovery_admin:'TempRecovery123!' \
  -XPOST "https://localhost:9200/_security/user/elastic/_password" \
  -H "Content-Type: application/json" \
  -d '{"password": "<新密码>"}'
\end{lstlisting}

使用完毕后删除临时用户：

\begin{lstlisting}[language=bash]
sudo /usr/share/elasticsearch/bin/elasticsearch-users userdel recovery_admin
\end{lstlisting}

\chapter{典型故障与排查}

\section{认证返回 401}

现象：使用安装时输出的密码或重置后的密码访问，均返回：

\begin{lstlisting}
security_exception: unable to authenticate user [elastic]
\end{lstlisting}

常见原因：
\begin{enumerate}[leftmargin=1.4em]
  \item 集群处于 RED 状态，\texttt{.security-7} 索引主分片未分配，密码写入失败；
  \item 首次安全初始化被中断，Elasticsearch 误判安全已完成并跳过初始化；
  \item 磁盘水位线超限，分片无法分配。
\end{enumerate}

排查顺序：先看集群健康，再看分片分配原因，最后确认磁盘空间。

\section{集群状态 RED 与 \texttt{.security-7} 不可用}

关键日志：

\begin{lstlisting}
UnavailableShardsException: [.security-7][0] primary shard is not active
\end{lstlisting}

说明安全索引主分片不可写，所有依赖安全索引的 API（包括密码重置）都会失败。

可查询未分配分片及其原因：

\begin{lstlisting}[language=bash]
curl "http://localhost:9200/_cat/shards?v&h=index,shard,prirep,state,unassigned.reason&s=state"
\end{lstlisting}

\section{磁盘水位线}

日志示例：

\begin{lstlisting}
high disk watermark [90%] exceeded on ... free: 3.1gb[5.5%],
shards will be relocated away from this node
\end{lstlisting}

Elasticsearch 默认高水位线为 90\%。可用空间低于该阈值时，节点拒绝分配新分片，包括 \texttt{.security-7}，进而导致集群 RED 与认证失败。

处理方式：释放磁盘空间，或将 \texttt{path.data} 迁移至空间充足的分区。仅在测试环境可临时关闭水位线检查：

\begin{lstlisting}[language=yaml]
cluster.routing.allocation.disk.threshold_enabled: false
\end{lstlisting}

\begin{warn}
关闭水位线检查只是绕过问题。生产环境必须保证磁盘余量充足。
\end{warn}

\section{数据目录迁移}

停止服务后整体移动数据目录，并修改 \texttt{path.data}：

\begin{lstlisting}[language=bash]
sudo systemctl stop elasticsearch
sudo mkdir -p /data/elasticsearch
sudo chown elasticsearch:elasticsearch /data/elasticsearch
sudo mv /var/lib/elasticsearch/* /data/elasticsearch/
\end{lstlisting}

\begin{lstlisting}[language=yaml]
path.data: /data/elasticsearch
\end{lstlisting}

启动后确认日志中不再出现针对旧路径的 watermark 警告，并确认集群状态恢复：

\begin{lstlisting}
Cluster health status changed from [YELLOW] to [GREEN]
(reason: [shards started [[.security-7][0], ...]])
\end{lstlisting}

\section{清理与重装}

无法通过 \texttt{apt} 时，使用 \texttt{dpkg} 彻底清除：

\begin{lstlisting}[language=bash]
sudo systemctl stop elasticsearch
sudo dpkg --purge elasticsearch
sudo rm -rf /etc/elasticsearch /var/lib/elasticsearch /var/log/elasticsearch
sudo rm -rf /usr/share/elasticsearch
sudo rm -f /lib/systemd/system/elasticsearch.service
sudo rm -f /etc/systemd/system/elasticsearch.service
sudo systemctl daemon-reload
sudo userdel elasticsearch
sudo groupdel elasticsearch
\end{lstlisting}

\begin{warn}
清理前确认没有需要保留的数据。重装后第一时间保存新生成的 \texttt{elastic} 初始密码。
\end{warn}

\chapter{常用检查命令}

\begin{center}
\begin{tabular}{@{}p{0.42\textwidth}p{0.5\textwidth}@{}}
\toprule
目的 & 命令 \\
\midrule
查看集群健康 & \texttt{GET /\_cluster/health?pretty} \\
查看红色索引 & \texttt{GET /\_cat/indices?v\&health=red} \\
查看未分配分片 & \texttt{GET /\_cat/shards?v\&status=unassigned} \\
分片分配原因 & \texttt{GET /\_cluster/allocation/explain} \\
查看安全索引 & \texttt{GET /\_cat/indices/.security*?v} \\
查看节点版本 & \texttt{GET /\_cat/nodes?v\&h=name,version} \\
\bottomrule
\end{tabular}
\end{center}

\begin{note}
URL 中的 \texttt{\&} 在 shell 中需加引号，否则会被当作后台运行符号截断命令。
\end{note}

\chapter{经验要点}

\begin{enumerate}[leftmargin=1.4em]
  \item 安装前先检查磁盘空间并规划 \texttt{path.data}，避免后续因水位线引发连锁故障。
  \item 单节点部署必须设置 \texttt{discovery.type: single-node}，并移除 \texttt{cluster.initial\_master\_nodes}。
  \item 不要通过关闭 \texttt{xpack.security.enabled} 来绕过认证问题，这会导致 \texttt{.security} 索引初始化失败，引发 RED 与后续认证循环。
  \item 认证 401 往往不是密码本身的问题，而是集群 RED 或安全索引不可写。
  \item 磁盘水位线是容易被忽略的根因，需优先确认 \texttt{path.data} 所在分区可用空间。
  \item file realm 应急用户是在安全索引损坏时恢复访问的可靠手段。
  \item 集群健康恢复为 GREEN 后，密码重置工具才能正常工作。
\end{enumerate}

\end{document}